\documentclass[11pt]{article}
\usepackage[margin=1in]{geometry}
\usepackage[T1]{fontenc}
\usepackage{lmodern}
\usepackage{microtype}
\usepackage{xcolor}
\usepackage{amsmath,amssymb}
\usepackage{enumitem}
\usepackage[colorlinks=true,linkcolor=blue!55!black,urlcolor=blue!55!black,citecolor=blue!55!black]{hyperref}

\definecolor{statusbg}{RGB}{239,246,255}
\setlist[itemize]{leftmargin=1.5em,itemsep=0.35em,topsep=0.35em}
\setlength{\parindent}{0pt}
\setlength{\parskip}{0.55em}
\setlength{\emergencystretch}{3em}

\begin{document}

\begin{center}
{\LARGE\bfseries Three convex Fuglede problems are answered in the literature\par}
\vspace{0.45em}
{\large arXiv:1602.08854 answered by arXiv:1904.12262\par}
\vspace{0.45em}
{\normalsize Run \texttt{fa\_banach\_001}, lane 2 \hfill 27 August 2026}
\end{center}

\vspace{0.6em}
\noindent\colorbox{statusbg}{%
\begin{minipage}{0.94\linewidth}
\textbf{Status: literature already answered.}
Theorem 1.1 of Lev and Matolcsi proves, in every dimension, that a spectral
convex body is a convex polytope and tiles face-to-face by translations along
a lattice. It therefore answers all three final-section problems below.
\end{minipage}}

\section*{Original source and exact questions}

Greenfeld and Lev, \emph{Fuglede's spectral set conjecture for convex
polytopes}, arXiv:1602.08854 \cite{GL}, Section 17, printed page 38, ask:

\begin{itemize}
\item \textbf{Problem 17.1.} If $\Omega\subset\mathbb{R}^d$, $d\geq 4$, is a
spectral convex polytope, prove that $\Omega$ tiles by translations.
\item \textbf{Problem 17.3.} Prove the same assertion when $\Omega$ is a
zonotope in $\mathbb{R}^d$, $d\geq 4$.
\item \textbf{Problem 17.4.} If a convex body $\Omega\subset\mathbb{R}^d$ is
spectral, prove that it must be a polytope.
\end{itemize}

Here ``spectral'' means that $L^2(\Omega)$ admits an orthogonal basis of
exponential functions. The first problem is the higher-dimensional
spectral-implies-tiling direction of Fuglede's conjecture restricted to convex
polytopes; the zonotope problem is its stated special case.

\section*{Decisive later theorem}

Lev and Matolcsi, \emph{The Fuglede conjecture for convex domains is true in
all dimensions}, arXiv:1904.12262 \cite{LM}, state:

\medskip
\noindent\textbf{Theorem 1.1.}
\emph{Let $\Omega$ be a convex body in $\mathbb{R}^d$. If $\Omega$ is a
spectral set, then $\Omega$ must be a convex polytope, and it tiles the space
face-to-face by translations along a lattice.}
\medskip

The polytope conclusion is exactly Problem 17.4. The tiling conclusion applies
to every spectral convex polytope in every dimension, which gives Problem
17.1, and hence its zonotope subcase Problem 17.3.

\newpage
\section*{Identification and provenance}

This is an explicit literature completion, not a new proof. The later paper
cites Greenfeld--Lev as the preceding three-dimensional polytope result and
describes the higher-dimensional polytope case as open before announcing its
all-dimensional theorem. One of its authors is also an author of the source.
The relationship therefore was known to the supporting authors rather than
being only an agent-discovered corollary.

The later theorem appeared in \emph{Acta Mathematica} 228 (2022), 385--420,
\url{https://doi.org/10.4310/ACTA.2022.v228.n2.a3}. Bounded searches by both arXiv identifiers,
the exact titles, and ``Fuglede convex domains all dimensions'' located this
direct answer.

\section*{Scope}

No part of Problems 17.1, 17.3, or 17.4 remains open. This identification says
nothing about unrestricted nonconvex spectral sets, where Fuglede's original
conjecture has a different status.

\begin{thebibliography}{9}
\bibitem{GL}
R. Greenfeld and N. Lev,
\emph{Fuglede's spectral set conjecture for convex polytopes},
arXiv:1602.08854; J. Anal. Math. 140 (2020), 409--464.

\bibitem{LM}
N. Lev and M. Matolcsi,
\emph{The Fuglede conjecture for convex domains is true in all dimensions},
arXiv:1904.12262; Acta Math. 228 (2022), 385--420.
\end{thebibliography}

\end{document}
